\begin{proof}
En otras palabras, la propiedad $\boldsymbol P$ para \emph{la fibra del punto
cerrado} de $\operatorname{Spec}(A)$ implica la misma propiedad para
\emph{todas las fibras} del morfismo $f$ (es decir, que $f$ es un
$\boldsymbol P$-\emph{morfismo} (7.3.1)). Procederemos en varias etapas:

I) \emph{Reducción al estudio de los anillos locales de la fibra
genérica.} --- Apliquemos el lema (7.3.16.2): basta comprobar que, para
toda $A$-álgebra \emph{finita} e \emph{íntegra} $A'$ con cuerpo de fracciones $K'$, si
se pone $B'=B\otimes_A A'$, todos los anillos locales de
$\operatorname{Spec}(B'\otimes_{A'}K')$ satisfacen la propiedad
$\boldsymbol R$. El anillo $A'$ (resp. $B'$) es semilocal
completo ($B'$ es una $B$-álgebra finita), y por tanto producto de
anillos locales completos; como $A'$ es íntegro por hipótesis, es \emph{local};
todo ideal maximal $\mathfrak n'$ de $B'$ está situado sobre el ideal
maximal $\mathfrak n$ de $B$ y, por tanto, sobre $\mathfrak m$; en
consecuencia, su imagen inversa en $A'$ es el ideal maximal
$\mathfrak m'$ de $A'$. Como todo anillo local de
$\operatorname{Spec}(B'\otimes_{A'}K')$ es un anillo local de
$\operatorname{Spec}(B')$, y por tanto de uno de los
$\operatorname{Spec}(B'_{\mathfrak n'})$ (sobre el punto genérico de
$A'$), se ve que bastará demostrar que los anillos locales de
$\operatorname{Spec}(B'_{\mathfrak n'}\otimes_{A'}K')$
poseen la propiedad $\boldsymbol R$. Ahora bien, $A'$ y
$B'_{\mathfrak n'}$ son anillos locales noetherianos completos; el
cuerpo residual de $B'_{\mathfrak n'}$ es una extensión finita del de
$B$, y por tanto también de $k$ y, \emph{a fortiori}, del cuerpo
residual $k'$ de $A'$. Esta observación y (2.1.4) muestran que $A'$ y
$B'_{\mathfrak n'}$ satisfacen las condiciones (i) y (ii) del
enunciado. Por otra parte, si $k''$ es una extensión finita de $k'$,
$k''$ es también una extensión finita de $k$, y los anillos locales de
$\operatorname{Spec}(B'_{\mathfrak n'}\otimes_{A'}k'')$
son también anillos locales de $\operatorname{Spec}(B\otimes_A k'')$;
por ello, la hipótesis de que se cumple
$\boldsymbol P(\operatorname{Spec}(B\otimes_A k),k)$ implica que la
propiedad
$\boldsymbol P(\operatorname{Spec}(B'_{\mathfrak n'}\otimes_{A'}k'),k')$
se cumple.

La cuestión se reduce así al caso en que, además, $A$ es \emph{íntegro}, con
cuerpo de fracciones $K$, y a demostrar que \emph{los anillos locales de}
$\operatorname{Spec}(B\otimes_A K)$ \emph{poseen la propiedad} $\boldsymbol R$.
